% from ega2-1-fr.tex
\oldpage[II]{75}

En particulier~:
\begin{enumerate}
  \item[1\textsuperscript{o}] Le corollaire
  (\hyperref[II.4.2.6-fr]{4.2.6}) s'applique lorsque $Y$ est un
  \emph{schéma local} (\hyperref[I.2.4.8-fr]{I, 2.4.8}). Soient $Y$ un
  préschéma quelconque, $y$ un point de $Y$,
  $Y'=\operatorname{Spec}(k(y))$~; la fibre $p^{-1}(y)$ de
  $\mathbf{P}(\mathcal{E})$ est, en vertu de
  (\hyperref[II.4.1.3.1-fr]{4.1.3.1}), identifiée à
  $\mathbf{P}(\mathcal{E}^y)$, où
  $\mathcal{E}^y=\mathcal{E}_y\otimes_{\mathcal{O}_y}k(y)
  =\mathcal{E}_y/\mathfrak{m}_y\mathcal{E}_y$ est considéré comme espace
  vectoriel sur $k(y)$. Plus généralement, si $K$ est une extension de
  $k(y)$, $p^{-1}(y)\otimes_{k(y)}K$ s'identifie à
  $\mathbf{P}(\mathcal{E}^y\otimes_{k(y)}K)$. Le cor.
  (\hyperref[II.4.2.6-fr]{4.2.6}) montre donc que l'ensemble des
  \emph{points géométriques de $\mathbf{P}(\mathcal{E})$ à valeurs dans
  l'extension $K$ de $k(y)$} (\hyperref[I.3.4.5-fr]{I, 3.4.5}), que l'on
  peut encore appeler \emph{fibre géométrique rationnelle sur $K$ de
  $\mathbf{P}(\mathcal{E})$ au-dessus du point $y$}, s'identifie à
  l'\emph{espace projectif} associé au \emph{dual} du $K$-espace vectoriel
  $\mathcal{E}^y\otimes_{k(y)}K$.
  \item[2\textsuperscript{o}] Supposons que $Y$ soit affine d'anneau $A$, et
  en outre que tout $\mathcal{O}_Y$-Module inversible soit trivial~; prenons
  en outre $\mathcal{E}=\mathcal{O}_Y^n$~; alors, dans
  (\hyperref[II.4.2.6-fr]{4.2.6}), $V$ s'identifie à $A^n$
  (\hyperref[I.1.3.8-fr]{I, 1.3.8}), $V^*$ à l'ensemble des systèmes
  $(f_i)_{1\leq i\leq n}$ d'éléments de $A$ engendrant l'idéal $A$~; deux
  tels systèmes définissent la même $Y$-section de
  $\mathbf{P}_Y^{n-1}=\mathbf{P}_A^{n-1}$, autrement dit le même point de
  $\mathbf{P}_A^{n-1}$ à valeurs dans $A$ si et seulement si l'un se déduit
  de l'autre par multiplication par un élément inversible de $A$.
\end{enumerate}

Ces propriétés justifient la terminologie de \og fibré projectif\fg{} pour
$\mathbf{P}(\mathcal{E})$. On notera que la définition que nous obtenons
ainsi de l'\og espace projectif\fg{} est en fait duale de la définition
classique~; cela nous est imposé par la nécessité de pouvoir définir
$\mathbf{P}(\mathcal{E})$ pour un $\mathcal{O}_Y$-Module quasi-cohérent
quelconque $\mathcal{E}$, non nécessairement localement libre.

\begin{remark}[4.2.7]
\label{II.4.2.7-fr}
Nous verrons au chapitre~V que si $Y$ est localement noethérien et connexe,
et si $\mathcal{E}$ est localement libre, alors, en posant
$P=\mathbf{P}(\mathcal{E})$, tout $\mathcal{O}_P$-Module inversible
$\mathcal{L}$ est isomorphe à un $\mathcal{O}_P$-Module de la forme
$\mathcal{L}'\otimes_{\mathcal{O}_Y}\mathcal{O}_P(m)$, où $\mathcal{L}'$ est
un $\mathcal{O}_Y$-Module inversible, bien déterminé à un isomorphisme près,
et $m$ un entier bien déterminé. En d'autres termes,
$\mathrm{H}^1(P,\mathcal{O}_P^*)$ est isomorphe à
$\mathbf{Z}\times\mathrm{H}^1(Y,\mathcal{O}_Y^*)$
(\hyperref[0.5.4.7-fr]{0, 5.4.7}). Nous verrons aussi
(\hyperref[III.2.1.14-fr]{III, 2.1.14}, compte tenu de
(\hyperref[0.5.4.10-fr]{0, 5.4.10})) que
$p_*(\mathcal{L}^{\otimes m})=0$ si $m<0$ et que
$p_*(\mathcal{L}^{\otimes m})$ est isomorphe à
$\mathcal{L}'\otimes_{\mathcal{O}_Y}
(\mathbf{S}_{\mathcal{O}_Y}(\mathcal{E}))_m$ si $m\geq0$. Si $\mathcal{F}$
est un $\mathcal{O}_Y$-Module quasi-cohérent, tout $Y$-morphisme
$\mathbf{P}(\mathcal{E})\to\mathbf{P}(\mathcal{F})$ est donc déterminé par
la donnée d'un $\mathcal{O}_Y$-Module inversible $\mathcal{L}'$, d'un entier
$m\geq0$ et d'un $\mathcal{O}_Y$-homomorphisme
$\psi:\mathcal{F}\to\mathcal{L}'\otimes_{\mathcal{O}_Y}
(\mathbf{S}_{\mathcal{O}_Y}(\mathcal{E}))_m$ tel que l'homomorphisme
correspondant $\psi^\sharp$ de $\mathcal{O}_{\mathbf{P}(\mathcal{F})}$-Modules
soit surjectif. Nous verrons aussi que si le $Y$-morphisme considéré est un
isomorphisme, alors $m=1$ et $\mathcal{F}$ est isomorphe à
$\mathcal{E}\otimes_{\mathcal{O}_Y}\mathcal{L}'$ (réciproque de
(\hyperref[II.4.1.4-fr]{4.1.4})). Ceci permettra de déterminer le faisceau
des germes d'automorphismes de $\mathbf{P}(\mathcal{E})$ comme quotient du
faisceau de groupes $\mathcal{A}ut(\mathcal{E})$ (localement isomorphe à
$\operatorname{GL}(n,\mathcal{O}_Y)$ si $\mathcal{E}$ est de rang $n$) par
$\mathcal{O}_Y^*$.
\end{remark}

\begin{env}[4.2.8]
\label{II.4.2.8-fr}
Conservant les notations de (\hyperref[II.4.2.1-fr]{4.2.1}), soit
$u:X'\to X$ un morphisme~; si le $Y$-morphisme $r:X\to P$ correspond à
l'homomorphisme $\varphi:q^*(\mathcal{E})\to\mathcal{L}$, le $Y$-morphisme
$r\circ u$ correspond à
$u^*(\varphi):u^*(q^*(\mathcal{E}))\to u^*(\mathcal{L})$, comme il résulte
aussitôt des définitions.
\end{env}

\begin{env}[4.2.9]
\label{II.4.2.9-fr}
Soient $\mathcal{E}$, $\mathcal{F}$ deux $\mathcal{O}_Y$-Modules
quasi-cohérents, $v:\mathcal{E}\to\mathcal{F}$ un homomorphisme surjectif,
$j=\mathbf{P}(v)$ l'immersion fermée correspondante
$\mathbf{P}(\mathcal{F})\to\mathbf{P}(\mathcal{E})$
(\hyperref[II.4.1.2-fr]{4.1.2}). Si le $Y$-morphisme
$r:X\to\mathbf{P}(\mathcal{F})$ correspond à l'homomorphisme
$\varphi:q^*(\mathcal{F})\to\mathcal{L}$, le

